\section{Conclusion}
\label{sec:conclusion}
We presented \capg{}, a color-based image immunization approach that offers an alternative to conventional additive perturbations through a smoothly constrained, per-image LUT. By optimizing color transformations under geometric constraints, \capg{} aims to disrupt unauthorized image editing while preserving image quality. Experiments demonstrate that, although additive defenses can achieve stronger disruption without purification, their effectiveness often degrades substantially after image purification. In contrast, \capg{} maintains stronger editing disruption and higher immunization success rates after both conventional image transformations and learned VAE-based purification. Moreover, its effectiveness also transfers to unseen editing models without re-optimization. Our findings open the door to exploring a broader design space of non-additive image immunization, moving beyond pixel-level perturbations toward protection mechanisms that remain effective under image transformations.

%We proposed \capg{}, which protects a portrait by a per-image 3D LUT that rotates colors smoothly about the gray axis within small bounds, so that the perturbation carries no spatial structure. On 306 portraits at matched distortion, this perturbation and the edit displacement it induces are preserved by purification: the edit displacement of \capg{} is 0.353 without purification, 0.364 on average over seven purifications, and 0.288 after VAE-Trans, against 0.112--0.275 and 0.023--0.143 for ten additive defenses. The ISR of \capg{} is 16.4\% without purification and 15.8\% after it, against 2.6--12.2\% for the five additive defenses evaluated after purification; without purification, it is not significantly higher than theirs.